% Source: Wald GR, physical PDF pp. 60--66
%         (printed pp. 47--53).
% Coverage: Section 3.4, Methods for Computing Curvature; equations (3.4.1)--(3.4.28).
% Figures/tables/footnotes: no figures, tables, or footnotes.
% Status: source-reviewed and compile-clean.
% Uncertainties: none.

\subsection{Methods for Computing Curvature}\label{sec:3.4}

In section 3.2 we defined the Riemann curvature tensor by proving existence of
a tensor field which satisfies Eq.~\eqref{eq:3.2.3} for all covector fields $\omega_a$. However,
this existence argument does not tell us directly how to calculate $R_{abc}{}^d$ given a metric
$g_{ab}$. Since the ability to calculate curvature is crucial for solving Einstein's equation
of general relativity, we devote this section to methods for calculating $R_{abc}{}^d$.

\medskip
\noindent\textbf{3.4a\quad Coordinate Component Method}

To calculate the curvature by the coordinate component method, we begin by
choosing a coordinate system. We express the derivative operator $\nabla_a$ in terms of the
ordinary derivative $\partial_a$ of this coordinate system and the Christoffel symbol $\Gamma^c{}_{ab}$, as
discussed in section 3.1. For a covector field $\omega_a$, we have
\begin{equation}
  \nabla_b\omega_c=\partial_b\omega_c-\Gamma^d{}_{bc}\omega_d,
  \tag{3.4.1}\label{eq:3.4.1}
\end{equation}
and thus,
\begin{equation}
  \begin{aligned}
    \nabla_a\nabla_b\omega_c
    ={}&\partial_a(\partial_b\omega_c-\Gamma^d{}_{bc}\omega_d)\\
    &-\Gamma^e{}_{ab}(\partial_e\omega_c-\Gamma^d{}_{ec}\omega_d)\\
    &-\Gamma^e{}_{ac}(\partial_b\omega_e-\Gamma^d{}_{be}\omega_d).
  \end{aligned}
  \tag{3.4.2}\label{eq:3.4.2}
\end{equation}
Hence, Eq.~\eqref{eq:3.2.3} may be expressed as
\begin{equation}
  R_{abc}{}^d\omega_d
  =\bigl[-2\partial_{[a}\Gamma^d{}_{b]c}
  +2\Gamma^e{}_{c[a}\Gamma^d{}_{b]e}\bigr]\omega_d,
  \tag{3.4.3}\label{eq:3.4.3}
\end{equation}
where we have used the commutativity of ordinary derivatives and the symmetry
of $\Gamma^c{}_{ab}$, Eq.~\eqref{eq:3.1.9}. Since Eq.~\eqref{eq:3.4.3} holds for all $\omega_d$, we may
\lq\lq cancel\rq\rq\ $\omega_d$ from both sides to obtain the desired formula for $R_{abc}{}^d$. Taking the components
of the tensors appearing in this equation in the coordinate basis associated with our
chart, we obtain the formula
\begin{equation}
  R_{\mu\nu\rho}{}^\sigma
  =\frac{\partial}{\partial x^\nu}\Gamma^\sigma{}_{\mu\rho}
  -\frac{\partial}{\partial x^\mu}\Gamma^\sigma{}_{\nu\rho}
  +\sum_\alpha\left(\Gamma^\alpha{}_{\mu\rho}\Gamma^\sigma{}_{\alpha\nu}
  -\Gamma^\alpha{}_{\nu\rho}\Gamma^\sigma{}_{\alpha\mu}\right),
  \tag{3.4.4}\label{eq:3.4.4}
\end{equation}
where we have used the definition of the ordinary derivative as the partial derivative
of components with respect to the coordinates.

Thus, to calculate $R_{abc}{}^d$ starting from $g_{ab}$, we first obtain the components, $g_{\mu\nu}$, of
the metric in our coordinate basis. We then calculate $\Gamma^\sigma{}_{\mu\nu}$ by Eq.~\eqref{eq:3.1.30} (or
by reading them off from the components of the geodesic equation~\eqref{eq:3.3.5} calculated
from the Lagrangian, Eq.~\eqref{eq:3.3.14}). Finally we calculate the components $R_{\mu\nu\rho}{}^\sigma$
via Eq.~\eqref{eq:3.4.4}. The Ricci tensor components then may be obtained by contracting Eq.~\eqref{eq:3.4.4}:
\begin{equation}
  \begin{aligned}
    R_{\mu\rho}
    &=\sum_\nu R_{\mu\nu\rho}{}^\nu\\
    &=\sum_\nu\frac{\partial}{\partial x^\nu}\Gamma^\nu{}_{\mu\rho}
    -\frac{\partial}{\partial x^\mu}\left(\sum_\nu\Gamma^\nu{}_{\nu\rho}\right)
    +\sum_{\alpha,\nu}\left(\Gamma^\alpha{}_{\mu\rho}\Gamma^\nu{}_{\alpha\nu}
    -\Gamma^\alpha{}_{\nu\rho}\Gamma^\nu{}_{\alpha\mu}\right).
  \end{aligned}
  \tag{3.4.5}\label{eq:3.4.5}
\end{equation}

We take this opportunity to point out several useful facts for calculations in
coordinate bases. We may write the coordinate basis components, $g_{\mu\nu}$, of the metric
as a matrix. The components, $g^{\mu\nu}$, of the inverse metric $g^{ab}$ will then be the inverse
of the matrix $(g_{\mu\nu})$. We define $g$ to be the determinant of $(g_{\mu\nu})$,
\begin{equation}
  g=\det(g_{\mu\nu}).
  \tag{3.4.6}\label{eq:3.4.6}
\end{equation}
Then, as discussed in appendix B, the natural volume element on the manifold $M$
induced by $g_{ab}$ is $\sqrt{|g|}\,dx^1\ldots dx^n$.

A simple formula may be derived for the contracted Christoffel symbol $\Gamma^a{}_{a\mu}$. From
Eq.~\eqref{eq:3.1.30} we have
\begin{equation}
  \Gamma^a{}_{a\mu}
  =\sum_\nu\Gamma^\nu{}_{\nu\mu}
  =\frac{1}{2}\sum_{\nu,\alpha}g^{\nu\alpha}
  \frac{\partial g_{\nu\alpha}}{\partial x^\mu}.
  \tag{3.4.7}\label{eq:3.4.7}
\end{equation}
But, using the formula for the inverse of a matrix, it is not difficult to show that
\begin{equation}
  \sum_{\nu,\alpha}g^{\nu\alpha}\frac{\partial g_{\nu\alpha}}{\partial x^\mu}
  =\frac{1}{g}\frac{\partial g}{\partial x^\mu}.
  \tag{3.4.8}\label{eq:3.4.8}
\end{equation}
Thus, we have
\begin{equation}
  \Gamma^a{}_{a\mu}
  =\frac{1}{2}\frac{1}{g}\frac{\partial g}{\partial x^\mu}
  =\frac{\partial}{\partial x^\mu}\ln\sqrt{|g|}.
  \tag{3.4.9}\label{eq:3.4.9}
\end{equation}
This is a useful formula since $\Gamma^a{}_{ab}$ appears in the expression for the Ricci tensor $R_{ab}$
as well as in the formula for the divergence of an arbitrary vector field, $T^a$, in terms
of its coordinate basis components. Indeed, using Eq.~\eqref{eq:3.4.9}, we have
\begin{equation}
  \begin{aligned}
    \nabla_aT^a
    &=\partial_aT^a+\Gamma^a{}_{ab}T^b\\
    &=\sum_\mu\frac{1}{\sqrt{|g|}}\frac{\partial}{\partial x^\mu}
    \bigl(\sqrt{|g|}T^\mu\bigr).
  \end{aligned}
  \tag{3.4.10}\label{eq:3.4.10}
\end{equation}

Coordinate basis methods have the advantage of providing a straightforward
mechanical procedure for calculating curvature. One is normally presented a metric
in terms of its components in a coordinate basis, so all one must do to obtain the
curvature is to perform the partial differentiations and summations given in equations
(3.1.30) and (3.4.4). On the other hand, even in the simplest cases, these straightforward computations are extremely laborious, and the \lq\lq nongeometrical\rq\rq\ nature of
the calculations often makes it difficult to detect algebraic errors.

\medskip
\noindent\textbf{3.4b\quad Orthonormal Basis (Tetrad) Methods}

While coordinate basis methods have the advantage of providing a straightforward
procedure for calculation of $\nabla_a$ and $R_{abc}{}^d$, for many purposes it is advantageous to use
an orthonormal basis in tensor calculations. A coordinate basis $\{\partial/\partial x^\mu\}$, of course,
is not orthonormal except for the trivial case of flat spacetime in Cartesian coordinates. Thus, we may wish to introduce a \lq\lq nonholonomic,\rq\rq\ i.e., noncoordinate,
orthonormal basis of smooth vector fields $\{(e_\mu)^a\}$, satisfying
\begin{equation}
  (e_\mu)^a(e_\nu)_a=\eta_{\mu\nu},
  \tag{3.4.11}\label{eq:3.4.11}
\end{equation}
where $\eta_{\mu\nu}=\diag(-1,\ldots,-1,1,\ldots,1)$. (Here the greek indices $\mu,\nu$ have the
range $1,\ldots,n$ and label the vector of the basis; the latin $a$ is the usual index of
the index notation designating that $e_\mu$ is a vector.) In four dimensions---and sometimes more generally---$\{(e_\mu)^a\}$ is called a \emph{tetrad}. Equation~\eqref{eq:3.4.11} implies the useful
relation
\begin{equation}
  \sum_{\mu,\nu}\eta^{\mu\nu}(e_\mu)^a(e_\nu)_b=\delta^a_b,
  \tag{3.4.12}\label{eq:3.4.12}
\end{equation}
where $\delta^a_b$ is the identity map on vectors and $\eta^{\mu\nu}$ is the inverse of $\eta_{\mu\nu}$ (so that
$\eta^{\mu\nu}=\eta_{\mu\nu}$). Equation~\eqref{eq:3.4.12} can be proven by verifying that it gives equality when
contracted with an arbitrary basis vector $(e_\sigma)^b$.

The calculation of curvature by tetrad methods has a sufficiently different appearance from the coordinate basis methods that it is worth pointing out explicitly three
key ingredients that must be used in determining the curvature of a metric: (1) The
derivative operator is compatible with the metric, $\nabla_ag_{bc}=0$. (2) The derivative
operator is torsion free (property 5 of section 3.1). (3) The Riemann tensor is related
to the derivative operator by Eq.~\eqref{eq:3.2.3}. In the coordinate basis methods,
ingredient (2) is expressed by Eq.~\eqref{eq:3.1.9}, ingredient (1) (given (2)) is expressed
by Eq.~\eqref{eq:3.1.29}, and ingredient (3) is expressed by Eq.~\eqref{eq:3.4.4}. As shown
below ingredients (1) and (2) are expressed in quite different ways when using an
orthonormal basis.

We begin by defining the \emph{connection 1-forms}, $\omega_{a\mu\nu}$, by
\begin{equation}
  \omega_{a\mu\nu}=(e_\mu)^b\nabla_a(e_\nu)_b.
  \tag{3.4.13}\label{eq:3.4.13}
\end{equation}
The components $\omega_{\lambda\mu\nu}$ of $\omega_{a\mu\nu}$ are called the \emph{Ricci rotation coefficients},
\begin{equation}
  \omega_{\lambda\mu\nu}=(e_\lambda)^a(e_\mu)^b\nabla_a(e_\nu)_b.
  \tag{3.4.14}\label{eq:3.4.14}
\end{equation}
The orthonormality of $\{(e_\mu)^a\}$ implies
\begin{equation}
  \begin{aligned}
    \omega_{a\mu\nu}
    &=(e_\mu)^b\nabla_a(e_\nu)_b\\
    &=-(e_\nu)^b\nabla_a(e_\mu)_b\\
    &=-\omega_{a\nu\mu},
  \end{aligned}
  \tag{3.4.15}\label{eq:3.4.15}
\end{equation}
where the compatibility condition $\nabla_ag_{bc}=0$ was used. Conversely, Eq.~\eqref{eq:3.4.15}
together with Eq.~\eqref{eq:3.4.11} implies $\nabla_ag_{bc}=0$. Thus, in the orthonormal
basis approach, ingredient (1) is expressed by the simple condition
\begin{equation}
  \omega_{a\mu\nu}=-\omega_{a\nu\mu}.
  \tag{3.4.16}\label{eq:3.4.16}
\end{equation}
(Note, in contrast, that the symmetry of $\Gamma^\alpha{}_{\mu\nu}$ stems from ingredient [2].) The
antisymmetry of the Ricci rotation coefficients as compared with the symmetry of the
Christoffel symbol indicates a significant potential advantage of this approach: There
are $n^2(n+1)/2$ ($=40$ when $n=4$) components $\Gamma^\lambda{}_{\mu\nu}$ but only $n^2(n-1)/2$ ($=24$
when $n=4$) components $\omega_{\lambda\mu\nu}$.

The Riemann tensor can be expressed in terms of the Ricci rotation coefficients
in the following manner. The components of $R_{abcd}$ in our orthonormal basis are
\begin{equation}
  \begin{aligned}
    R_{\rho\sigma\mu\nu}
    &=R_{abcd}(e_\rho)^a(e_\sigma)^b(e_\mu)^c(e_\nu)^d\\
    &=(e_\rho)^a(e_\sigma)^b(e_\mu)^c
    (\nabla_a\nabla_b-\nabla_b\nabla_a)(e_\nu)_c.
  \end{aligned}
  \tag{3.4.17}\label{eq:3.4.17}
\end{equation}
However, we have
\begin{equation}
  (e_\mu)^c\nabla_a\nabla_b(e_\nu)_c
  =\nabla_a\bigl\{(e_\mu)^c\nabla_b(e_\nu)_c\bigr\}
  -\bigl[\nabla_a(e_\mu)^c\bigr]\bigl[\nabla_b(e_\nu)_c\bigr].
  \tag{3.4.18}\label{eq:3.4.18}
\end{equation}
Furthermore, we have
\begin{equation}
  \begin{aligned}
    \bigl[\nabla_a(e_\mu)^c\bigr]\bigl[\nabla_b(e_\nu)_c\bigr]
    &={}\bigl[\nabla_a(e_\mu)^f\bigr]\delta^c_f
    \bigl[\nabla_b(e_\nu)_c\bigr]\\
    &={}\sum_{\alpha,\beta}\eta^{\alpha\beta}
    \bigl[\nabla_a(e_\mu)^f\bigr](e_\alpha)^c(e_\beta)_f
    \bigl[\nabla_b(e_\nu)_c\bigr].
  \end{aligned}
  \tag{3.4.19}\label{eq:3.4.19}
\end{equation}
Thus, from the definition of the connection one forms, Eq.~\eqref{eq:3.4.13}, we obtain
\begin{equation}
  \begin{aligned}
    R_{\rho\sigma\mu\nu}
    ={}&(e_\rho)^a(e_\sigma)^b\biggl\{
    \nabla_a\omega_{b\mu\nu}-\nabla_b\omega_{a\mu\nu}\\
    &\qquad-\sum_{\alpha,\beta}\eta^{\alpha\beta}
    \bigl[\omega_{a\beta\mu}\omega_{b\alpha\nu}
    -\omega_{b\beta\mu}\omega_{a\alpha\nu}\bigr]\biggr\}.
  \end{aligned}
  \tag{3.4.20}\label{eq:3.4.20}
\end{equation}
We may rewrite Eq.~\eqref{eq:3.4.20} in terms of the Ricci rotation coefficients as
\begin{equation}
  \begin{aligned}
    R_{\rho\sigma\mu\nu}
    ={}&(e_\rho)^a\nabla_a\omega_{\sigma\mu\nu}
    -(e_\sigma)^a\nabla_a\omega_{\rho\mu\nu}\\
    &-\sum_{\alpha,\beta}\eta^{\alpha\beta}\bigl\{
    \omega_{\rho\beta\mu}\omega_{\sigma\alpha\nu}
    -\omega_{\sigma\beta\mu}\omega_{\rho\alpha\nu}\\
    &\qquad\qquad+\omega_{\rho\beta\sigma}\omega_{\alpha\mu\nu}
    -\omega_{\sigma\beta\rho}\omega_{\alpha\mu\nu}\bigr\},
  \end{aligned}
  \tag{3.4.21}\label{eq:3.4.21}
\end{equation}
where the last two terms compensate for taking the components of $\omega_{a\mu\nu}$ inside the
derivative in the first two terms. Since $\omega_{\sigma\mu\nu}$ is a scalar, the derivative $\nabla_a$ in Eq.~\eqref{eq:3.4.21}
may, of course, be replaced by an ordinary derivative $\partial_a$. Thus, Eq.~\eqref{eq:3.4.21}
tells us how to compute the curvature tensor in terms of $\omega_{\sigma\mu\nu}$. The components of the Ricci tensor may then be computed via the formula
\begin{equation}
  R_{\rho\mu}=\sum_{\sigma,\nu}\eta^{\sigma\nu}R_{\rho\sigma\mu\nu}.
  \tag{3.4.22}\label{eq:3.4.22}
\end{equation}

Equation~\eqref{eq:3.4.20} or Eq.~\eqref{eq:3.4.21} fully expresses ingredient (3) in the calculation of curvature. Thus, only ingredient (2)---the torsion-free condition on $\nabla_a$---remains to be expressed in this approach. There are two different procedures which
may be employed to do this. The first begins by noting that the torsion-free condition
allowed us to express the commutator of two vector fields in terms of the derivative
operator $\nabla_a$ via Eq.~\eqref{eq:3.1.2}. Conversely, if Eq.~\eqref{eq:3.1.2} holds for all vector
fields in a basis, it implies the torsion-free property of $\nabla_a$. Thus, we may express
ingredient (2) by the \emph{commutation relations} of the basis vector fields,
\begin{equation}
  \begin{aligned}
    (e_\sigma)_a[e_\mu,e_\nu]^a
    &=(e_\sigma)_a\bigl\{(e_\mu)^b\nabla_b(e_\nu)^a
    -(e_\nu)^b\nabla_b(e_\mu)^a\bigr\}\\
    &=\omega_{\mu\sigma\nu}-\omega_{\nu\sigma\mu}.
  \end{aligned}
  \tag{3.4.23}\label{eq:3.4.23}
\end{equation}
This yields the $n^2(n-1)/2$ equations needed to solve for $\omega_{\sigma\mu\nu}$ (see problem 8).

An alternative procedure is to note that from the definition of the connection
one-forms, it follows that
\begin{equation}
  \nabla_{[a}(e_\sigma)_{b]}
  =\sum_{\mu,\nu}\eta^{\mu\nu}(e_\mu)_{[a}\omega_{b]\sigma\nu},
  \tag{3.4.24}\label{eq:3.4.24}
\end{equation}
as can be verified by contracting Eq.~\eqref{eq:3.4.24} with arbitrary basis vectors
$(e_\rho)^a(e_\lambda)^b$. However, the torsion-free condition implies that the antisymmetrized
derivative of a one-form is independent of derivative operator, and thus may be
replaced by an ordinary derivative $\partial_a$,
\begin{equation}
  \partial_{[a}(e_\sigma)_{b]}
  =\sum_{\mu,\nu}\eta^{\mu\nu}(e_\mu)_{[a}\omega_{b]\sigma\nu}.
  \tag{3.4.25}\label{eq:3.4.25}
\end{equation}
Conversely, the validity of Eq.~\eqref{eq:3.4.25} for all basis vectors implies that $\nabla_a$ is
torsion free. Thus, Eq.~\eqref{eq:3.4.25} is an alternate expression of ingredient (2).

Equations~\eqref{eq:3.4.20} and~\eqref{eq:3.4.25} can be expressed more elegantly in the notation
of differential forms (see appendix B). We drop the covector index $a$ on $\omega_{a\mu\nu}$
and $(e_\mu)_a$ and use boldface letters to designate forms. For convenience, we also raise and
lower greek indices with $\eta^{\mu\nu}$ and $\eta_{\mu\nu}$, e.g., we have
\begin{equation}
  \omega_\nu{}^\mu=\sum_\sigma\eta^{\mu\sigma}\omega_{\nu\sigma}.
  \tag{3.4.26}\label{eq:3.4.26}
\end{equation}
Then Eq.~\eqref{eq:3.4.25} becomes, in differential forms notation,
\begin{equation}
  d\boldsymbol{e}_\sigma
  =\sum_\mu\boldsymbol{e}_\mu\wedge\boldsymbol{\omega}_\sigma{}^\mu.
  \tag{3.4.27}\label{eq:3.4.27}
\end{equation}
Furthermore, Eq.~\eqref{eq:3.4.20} can be written as
\begin{equation}
  \boldsymbol{R}_\mu{}^\nu
  =d\boldsymbol{\omega}_\mu{}^\nu
  +\sum_\alpha\boldsymbol{\omega}_\mu{}^\alpha\wedge\boldsymbol{\omega}_\alpha{}^\nu,
  \tag{3.4.28}\label{eq:3.4.28}
\end{equation}
where $\boldsymbol{R}_{\mu\nu}$ is the two-form corresponding to $R_{ab\mu\nu}$. Equations~\eqref{eq:3.4.27}
and~\eqref{eq:3.4.28} are sometimes referred to as the \emph{equations of structure}.

Thus, we have basically two procedures for calculating curvature with orthonormal bases. The first is the straightforward approach by which one solves for the
Ricci rotation coefficients from the commutation relations, Eq.~\eqref{eq:3.4.23}, and
substitutes in Eq.~\eqref{eq:3.4.21}. The second approach---first discussed in the context
of calculations in general relativity by Misner (1963)---uses Eq.~\eqref{eq:3.4.27}
($=$ Eq.~\eqref{eq:3.4.25}) to determine the connection one-forms, from which the curvature may be
calculated by Eq.~\eqref{eq:3.4.28} ($=$ Eq.~\eqref{eq:3.4.20}). The major advantage of the second
approach is that it is relatively easy to calculate the left side of Eq.~\eqref{eq:3.4.27}. In
simple cases it is possible to guess the solution for $\boldsymbol{\omega}_\mu{}^\nu$ without further calculations.
In this way, a great deal of labor (such as that required to solve Eq.~\eqref{eq:3.4.23} or, worse
yet, the labor required to calculate $\Gamma^\alpha{}_{\mu\nu}$) can be avoided. We will illustrate this in
chapter 6. Note that in both approaches, one will still introduce a coordinate system
when it comes time to do the calculations. The aim of tetrad methods is not to avoid
the introduction of a coordinate system but rather to use a more convenient basis for
the expansion of tensors than that provided by the coordinate basis.

A variant of the \lq\lq straightforward approach\rq\rq---motivated by spinor methods (see
chapter 13) and specially adapted to the case of a Lorentz metric in four
dimensions---was introduced by Newman and Penrose (1962). Instead of choosing
an orthonormal basis, one chooses a \lq\lq null tetrad\rq\rq\ consisting of two null vectors
$l^a$ and $n^a$ with $l^an_a=-1$; one completes the basis by picking two orthonormal
spacelike vectors $x^a$ and $y^a$ orthogonal to $l^a$ and $n^a$ and defines the complex vector
$m^a=(x^a+\ii y^a)/\sqrt{2}$. The analogs of the 24 real Ricci rotation coefficients of an
orthonormal tetrad are 12 complex quantities called \emph{spin coefficients}. They are
related to the null tetrad by the analog of Eq.~\eqref{eq:3.4.13} and the curvature is given
in terms of them by the analog of Eq.~\eqref{eq:3.4.21}. The precise equations are given
by Newman and Penrose (1962), Pirani (1965), and others. The Newman-Penrose
approach has proven very useful for many applications in general relativity.

Which method---coordinate basis, orthonormal basis, or null tetrad---should one
use to do calculations? There is no universal answer. In many simple cases the
second orthonormal basis approach---using Eqs.~\eqref{eq:3.4.27} and~\eqref{eq:3.4.28}---saves
considerable labor over the other approaches, provided one can solve Eq.~\eqref{eq:3.4.27}
for the connection one-forms by inspection. In situations where a null vector
plays a geometrically distinguished role---specifically in algebraically special spacetimes (see chapter 7)---the Newman-Penrose approach is likely to be the most useful.
Furthermore, in some cases where there is a great deal of symmetry, it may be
possible to adapt the coordinate system to the symmetries and make the coordinate
basis method the preferred approach. Finally in general situations---without symmetries or geometrically preferred vectors to serve as natural basis vectors---all
methods are likely to be about equally laborious.

Finally, we note that in this section we have presented the coordinate and tetrad
methods in the context of calculating the Riemann tensor, given the metric. However, the equations obtained, of course, are equally applicable in the much more
frequently encountered situation where (part of) the curvature tensor is given (e.g.,
via Einstein's equation, as discussed in the next chapter) and we wish to find the
metric. Furthermore, even in cases where we do not wish to calculate the curvature
or solve for the metric, the above equations often directly provide useful relations
between the metric and its curvature, particularly in cases where the coordinate
system or tetrad is well adapted to the properties of the spacetimes under consideration.
